\section{Locate both quantizers on one grid}\label{ball-x-quantizers}
\begin{figure}[htbp]
  \centering
  \begin{tikzpicture}[alt={Two eight-column eight-row coordinate grids show the same raw position x 320, y 448. Floor highlights column one, row one, ID nine. Nearest highlights column one, row two, ID seventeen. All other cells are unselected.},x=6mm,y=6mm]
    \foreach \panel/\selected in {0/1,12/2} {
        \foreach \r in {0,1,2,3,4,5,6,7} {
            \foreach \c in {0,1,2,3,4,5,6,7} {\draw[gray] (\panel+\c,-\r) rectangle +(1,-1);}
            \node[font=\scriptsize] at (\panel-0.6,-\r-0.5) {\r};
          }
        \foreach \c in {0,1,2,3,4,5,6,7} {\node[font=\scriptsize] at (\panel+\c+0.5,0.5) {\c};}
        \fill (\panel+1,-\selected) rectangle +(1,-1);
      }
    \node[above] at (4,1) {Floor};\node[above] at (16,1) {Nearest};
    \node[text=white,font=\scriptsize] at (1.5,-1.5) {9};
    \node[text=white,font=\scriptsize] at (13.5,-2.5) {17};
    \node[below] at (4,-8.4) {Cell \((1,1)\)};
    \node[below] at (16,-8.4) {Cell \((1,2)\)};
  \end{tikzpicture}
  \caption{Coordinate diagrams highlight the cells chosen from one unchanged raw position. They illustrate geometry; actual IDE display captures remain separate pending figures.}
\end{figure}
Both renderers read raw \((320,448)\). Floor drops the fractions, while nearest
adds half a cell before dropping them. The main state words must be identical in
the two runs. The selected ID and frame can differ without changing the velocity
or the next raw position.
